% !TeX root = ../slides.tex
% =============================================================================
% Conclusões — um slide. Responde às perguntas (A) e (B) da Metodologia com os
% números dos slides de resultados. Cada número leva a taxa de poda em que vale.
% =============================================================================
\newcommand{\pergunta}[2]{%
  {\fontsize{8.5}{10.5}\selectfont\color{destaque}\textbf{Pergunta #1}\par}\vspace{0.8mm}
  {\fontsize{9}{11.5}\selectfont\color{tintaSec}\raggedright #2\par}\vspace{2.2mm}}
\newcommand{\resposta}[1]{{\fontsize{10.5}{13}\selectfont\bfseries\raggedright #1\par}\vspace{2.2mm}}

\begin{frame}
\framesubtitle{Conclusões}
\frametitle{A poda estruturada reduz a energia, e a análise de sensibilidade reduz
a perda de qualidade}
\vspace{1mm}
\begin{columns}[T,onlytextwidth]
\begin{column}{0.48\textwidth}
\pergunta{A}{A redução de parâmetros se converte em redução real de tempo e
energia?}
\resposta{Só quando a poda remove estruturas.}
\begin{itemize}
  \setlength{\itemsep}{2mm}
  \item Poda estruturada a 90\%: 61\% menos parâmetros, \textbf{59\% menos energia} e
        56\% menos tempo.
  \item Poda por magnitude: 61\% menos parâmetros não nulos, sem redução detectável
        de energia.
\end{itemize}
\end{column}
\begin{column}{0.48\textwidth}
\pergunta{B}{A análise de sensibilidade por camada melhora o compromisso entre
compressão e qualidade?}
\resposta{Sim, nas taxas de poda de 10\% a 50\%.}
\begin{itemize}
  \setlength{\itemsep}{2mm}
  \item Poda estruturada: \ppl{} \textbf{2,5 vezes menor} na média, até 4 vezes
        a 20\%, com a mesma energia.
  \item Poda por magnitude: menor \ppl{} em 7 de 9 taxas.
  \item Após uma época de \ft, a 30\%: \ppl{} 1,87 vez menor.
\end{itemize}
\end{column}
\end{columns}
\end{frame}
